% !TEX root = ../../main_power.tex
% ---------------------------------------------------------------------
%  Software-Massnahmen und Problembereiche (Teil 4)
% ---------------------------------------------------------------------
\section{Software}
\label{sec:pwrsw}

\begin{formel}[Das Leitmotiv]
\T{Schnell und kurz arbeiten, lange schlafen.} Weil der Low-Power-Modus
um Grössenordnungen sparsamer ist als der Betrieb, zahlt sich jede
Verkürzung der Wachphase doppelt aus -- selbst wenn sie mehr
Spitzenleistung kostet.
\end{formel}

\begin{ztab}{@{}lX@{}}
Sleep & so früh wie möglich hinein, so spät wie möglich heraus \\
\T{Gear Shifting} & Takt dynamisch an die Last anpassen, analog dem
  Blutdruck-/Herzschlagsystem: schnell nur, wenn nötig \\
Delays & \emph{nicht} aktiv pollen \ra Interrupt \\
Rechnen & Look-up-Tabellen statt Rechnung, FPU oder Coprozessor statt
  Software-Emulation \\
Peripherie & interne Einheiten abschalten; Resultat per Interrupt
  abholen statt auf Fertigmeldung zu warten \\
Interrupts & zum Aufwachen -- welche Quelle weckt, ist
  hardwareabhängig \\
Code-Dichte & weniger Instruktionen, Compiler-Optimierungen, Makros \\
\end{ztab}

\subsection{Compiler}
Es gibt keine Option „optimiere auf Low Power“. Die beiden vorhandenen
Richtungen wirken gegenläufig:
\begin{ztab}{@{}lX@{}}
auf Code-\emph{Grösse} & langsamerer Code, aber weniger davon -- besser
  für Cache und Flash \\
auf \emph{Geschwindigkeit} & grösserer Code, aber die Arbeit ist früher
  fertig \ra früher im Low-Power-Modus \\
\end{ztab}
\merkzeile Welche Richtung gewinnt, lässt sich nicht vorhersagen, nur
messen. Das ist genau die Art Aussage, die im Laborjournal belegt
gehört.

\begin{bsp}[Energieeffizienz von Programmiersprachen]
Normiert auf C $=$ 1.00 (Studie „Energy Efficiency across Programming
Languages“):
\begin{tabular}{@{}lr@{\hskip8pt}lr@{}}
C & 1.00 & C\# & 3.14 \\
Rust & 1.03 & Go & 3.23 \\
C++ & 1.34 & JavaScript & 4.45 \\
Ada & 1.70 & TypeScript & 21.50 \\
Java & 1.98 & \textbf{Python} & \textbf{75.88} \\
\end{tabular}\\[2pt]
Die Sprachwahl ist damit eine Grössenordnung wert -- mehr als die
meisten Mikrooptimierungen im Code.
\end{bsp}

\subsection{Problembereiche}
\label{falle:pwr}
\begin{falle}[Woran Low-Power-Designs scheitern]
\begin{description}
\item[Aufwachen] dauert -- ein 32-kHz-Quarz braucht lange zum
  Anschwingen. Trade-off: zu häufiges Aufwachen frisst den Gewinn des
  Schlafens auf.
\item[Clock] im Sparmodus eingeschränkte Funktionen; beim Gear Shifting
  ziehen die Prescaler aller Peripherieeinheiten (UART!) mit.
\item[Speisung] interne Schutzdioden, Stützkondensatoren, Einkopplung
  von Signalpegeln von aussen.
\item[Messung] Innenwiderstand des Multimeters (Spannungsabfall oder
  Kurzschluss) und der \emph{Dynamikbereich}: \textmu A im Schlaf und
  mA im Betrieb im selben Messbereich.
\item[UART] Break-Zustand braucht Pull-up; vor dem Schlafen warten, bis
  alles gesendet ist.
\item[Pulls] I\textsuperscript{2}C hochohmig auslegen -- aber nicht zu
  hochohmig, sonst leiden die Flanken.
\item[Pins] intern nicht verdrahtete Pins nicht treiben; Eingang mit
  internem Pull-up je nach Hersteller; Debug-Pins beachten.
\item[Batterie] Selbstentladung, und die Spannung ist direkt nach dem
  Einschalten höher als im eingeschwungenen Zustand.
\item[Dioden] Z-Dioden leiten unterhalb der Z-Spannung bereits im
  \textmu A-Bereich; dasselbe bei in Sperrrichtung geschalteten
  Schottky-Dioden.
\item[Reset] im Schlafmodus empfindlich gegen elektrostatische
  Aufladung \ra Pull-ups und Filter.
\end{description}
\end{falle}
